\item Un rayo de luz entra a la atmósfera de un planeta y desciende verticalmente a la superficie que está a una distancia $h$ hacia abajo. El índice de refracción donde la luz entra a la atmósfera es de $1.00$ y aumenta linealmente con la distancia a la superficie, donde tiene un valor de $n$.
\begin{enumerate}
    \item ¿Cuánto tarda el rayo en recorrer esta trayectoria?
    \item ¿En qué porcentaje es mayor el intervalo de tiempo que el necesario en la ausencia de una atmósfera?
\end{enumerate}